% GENERATED FILE. Edit canonical lessons, then run npm run generate:books.
% canonical-source-hash: fnv1a64:b3cd223d5f0dc933
% canonical-lessons: ZH-C45-tian, ZH-W45-tian, ZH-C45-jintian, ZH-W45-jin, ZH-C45-zuotian, ZH-C45-xianzai, ZH-C45-mingtian

\chapter{Today and Tomorrow}
\label{ch:zh-today-and-tomorrow}
\hlchaptermodality{45}

\begin{chapteropening}
\textbf{By the end of this chapter:} \emph{I can say today, yesterday, tomorrow and now.}

Complete the last lesson of chapter 45: I can say today, yesterday, tomorrow and now.
\end{chapteropening}

\section[\texorpdfstring{\zh{天} (tiān)}{天 (tiān)}]{\zh{天} (tiān) — day; sky}

\label{lesson:ZH-C45-tian}



\begin{quote}
Before the new one: say the Mandarin for to sit, then the Mandarin for taxi.
\end{quote}

\subsection*{You'll want to know: \zh{天}}
\textbf{\zh{天}} — \emph{tiān} — \textquotedblleft{}day; sky\textquotedblright{}.

\begin{quote}
\textbf{\zh{那天我们去大学了}} \emph{nà tiān wǒmen qù dàxué le} — That day we went to the university.
\end{quote}

The new character, in pieces: \zh{天} — \zh{一} above \zh{大}: one more bar over the big piece. The next lesson puts it in your hand.

\begin{cousinweb}
\textbf{\zh{天}} (\emph{tiān}) is inherited from Middle Chinese \emph{then} (Guangyun, in Baxter's transcription), glossed 'the highest, with nothing above it', the sky. Its use for a day is not given there.
\end{cousinweb}

\begin{grammarlens}[title={What you've built}]
One more word for times and partings.
\end{grammarlens}

\subsection*{Your turn}
\begin{itemize}
\raggedright
  \item \emph{Say it:} \emph{tiān}
  \item \emph{Say it:} \emph{tiān}, once more
  \item \emph{Say it:} \emph{nà tiān wǒmen qù dàxué le} — That day we went to the university.
  \item \emph{Recall:} say the Mandarin for to sit, then the Mandarin for taxi, then say \emph{tiān} again
\end{itemize}

\begin{tcolorbox}[breakable,colback=teal!4,colframe=teal!35!black,title={Before you move on}]
What does \zh{天} mean? (\textquotedblleft{}Day; sky\textquotedblright{}.) Say it once more.
\end{tcolorbox}
\section[\texorpdfstring{\zh{天} (tiān)}{天 (tiān)}]{\zh{天} — one more bar over \zh{大}}

\label{lesson:ZH-W45-tian}



\begin{quote}
The word first, then the character inside it.

\begin{itemize}
\raggedright
  \item \emph{Recall:} say \emph{tiān} — \textbf{R1}, one lesson back, before it decays
\end{itemize}

\begin{quote}
\zh{一} \zh{大}
\end{quote}

\emph{Read it:} each one aloud, with nothing said first — \emph{yī}, \emph{dà}
\end{quote}

\subsection*{Script — one more bar over \zh{大}}
\hlblockfigure{\detokenize{figures/ZH-W45-tian-filmstrip.pdf}}{How \zh{天} is written, stroke by stroke}

\begin{quote}
\zh{天}
\end{quote}

\textbf{Two pieces your hand knows, top to bottom: \zh{一}, then \zh{大}.}

\textbf{\zh{一}, one stroke}, a short horizontal at the top.

\textbf{\zh{大}, three strokes}: a longer horizontal, a left-falling stroke from above it, and a right-falling stroke.

1 + 3 = four. Keep the top bar shorter than the one under it.

\subsection*{Your turn}
\begin{itemize}
\raggedright
  \item \emph{Write it:} \zh{天}, all four, with the model in front of you
  \item \emph{Cover:} write it again from memory, and count the pieces, not the strokes
\end{itemize}

\begin{tcolorbox}[breakable,colback=teal!4,colframe=teal!35!black,title={Before you move on}]
Which two characters you know make \zh{天}? (\textbf{\zh{一} and \zh{大}}.) Which bar is longer? (\textbf{The lower one}.) How many strokes? (\textbf{Four}.)

Source: \href{https://raw.githubusercontent.com/chanind/hanzi-writer-data/68d10a4b21150cae5e1ebbd223eed289cf32d90c/data/\%E5\%A4\%A9.json}{Hanzi Writer Data \zh{天}, pinned stroke paths and medians}.
\end{tcolorbox}
\section[\texorpdfstring{\zh{今天} (jīntiān)}{今天 (jīntiān)}]{\zh{今天} (jīntiān) — today}

\label{lesson:ZH-C45-jintian}



\begin{quote}
Before the new one: say the Mandarin for taxi, then the Mandarin for day.
\end{quote}

\subsection*{You'll want to know: \zh{今天}}
\textbf{\zh{今天}} — \emph{jīntiān} — \textquotedblleft{}today\textquotedblright{}.

\begin{quote}
\textbf{\zh{你们今天去吗}} \emph{nǐmen jīntiān qù ma} — Are you going today?
\end{quote}

The new character, in pieces: \zh{今} — \zh{人} as a roof, a dot under it, and a stroke that runs across and turns down. The next lesson puts it in your hand.

\begin{cousinweb}
\textbf{\zh{今天}} (\emph{jīntiān}) is built from \zh{今} 'now' and \zh{天} 'day', both inherited from Middle Chinese, \emph{kim} and \emph{then} (Guangyun, in Baxter's transcription); the dictionary glosses \emph{kim} as the opposite of olden times.
\end{cousinweb}

\begin{grammarlens}[title={What you've built}]
One more word for times and partings.
\end{grammarlens}

\subsection*{Your turn}
\begin{itemize}
\raggedright
  \item \emph{Say it:} \emph{jīntiān}
  \item \emph{Say it:} \emph{jīntiān}, once more
  \item \emph{Say it:} \emph{nǐmen jīntiān qù ma} — Are you going today?
  \item \emph{Recall:} say the Mandarin for taxi, then the Mandarin for day, then say \emph{jīntiān} again
\end{itemize}

\begin{tcolorbox}[breakable,colback=teal!4,colframe=teal!35!black,title={Before you move on}]
What does \zh{今天} mean? (\textquotedblleft{}Today\textquotedblright{}.) Say it once more.
\end{tcolorbox}
\section[\texorpdfstring{\zh{今} (jīn)}{今 (jīn)}]{\zh{今} — a roof, a dot and a turn}

\label{lesson:ZH-W45-jin}



\begin{quote}
The word first, then the character inside it.

\begin{itemize}
\raggedright
  \item \emph{Recall:} say \emph{jīntiān} — \textbf{R1}, one lesson back, before it decays
\end{itemize}

\begin{quote}
\zh{人} \zh{个}
\end{quote}

\emph{Read it:} each one aloud, with nothing said first — \emph{rén}, \emph{gè}
\end{quote}

\subsection*{Script — a roof, a dot and a turn}
\hlblockfigure{\detokenize{figures/ZH-W45-jin-filmstrip.pdf}}{How \zh{今} is written, stroke by stroke}

\begin{quote}
\zh{今}
\end{quote}

\textbf{Top to bottom.}

\textbf{The roof, \zh{人}, two strokes}, spread wide: a left-falling stroke, then a right-falling stroke, as in \zh{个}.

\textbf{A dot} in the middle, under the roof.

\textbf{A stroke that runs across and turns down to the left}, under the dot.

2 + 1 + 1 = four.

\subsection*{Your turn}
\begin{itemize}
\raggedright
  \item \emph{Write it:} \zh{今}, all four, with the model in front of you
  \item \emph{Cover:} write it again from memory, and count the pieces, not the strokes
\end{itemize}

\begin{tcolorbox}[breakable,colback=teal!4,colframe=teal!35!black,title={Before you move on}]
Which character you know is the roof of \zh{今}? (\textbf{\zh{人}}.) Which other character has the same roof? (\textbf{\zh{个}}.) What comes right after the roof? (\textbf{A dot}.)

Source: \href{https://raw.githubusercontent.com/chanind/hanzi-writer-data/68d10a4b21150cae5e1ebbd223eed289cf32d90c/data/\%E4\%BB\%8A.json}{Hanzi Writer Data \zh{今}, pinned stroke paths and medians}.
\end{tcolorbox}
\section[zuótiān]{zuótiān — yesterday}

\label{lesson:ZH-C45-zuotian}



\begin{quote}
Before the new one: say the Mandarin for day, then the Mandarin for today.
\end{quote}

\subsection*{You'll want to know: zuótiān}
\emph{zuótiān} — \textquotedblleft{}yesterday\textquotedblright{}.

\begin{quote}
\emph{zuótiān wǒ hěn lèi} — I was really tired yesterday.
\end{quote}

Its characters come later; for now it is a word for the ear.

\begin{cousinweb}
\emph{zuótiān} is built from zuó 'the day before' and \zh{天} 'day', both inherited from Middle Chinese, \emph{dzak} and \emph{then} (Guangyun, in Baxter's transcription); the dictionary glosses \emph{dzak} as the day one night back.
\end{cousinweb}

\begin{grammarlens}[title={What you've built}]
One more word for times and partings.
\end{grammarlens}

\subsection*{Your turn}
\begin{itemize}
\raggedright
  \item \emph{Say it:} \emph{zuótiān}
  \item \emph{Say it:} \emph{zuótiān}, once more
  \item \emph{Say it:} \emph{zuótiān wǒ hěn lèi} — I was really tired yesterday.
  \item \emph{Recall:} say the Mandarin for day, then the Mandarin for today, then say \emph{zuótiān} again
\end{itemize}

\begin{tcolorbox}[breakable,colback=teal!4,colframe=teal!35!black,title={Before you move on}]
What does zuótiān mean? (\textquotedblleft{}Yesterday\textquotedblright{}.) Say it once more.
\end{tcolorbox}
\section[xiànzài]{xiànzài — now}

\label{lesson:ZH-C45-xianzai}



\begin{quote}
Before the new one: say the Mandarin for today, then the Mandarin for yesterday.
\end{quote}

\subsection*{You'll want to know: xiànzài}
\emph{xiànzài} — \textquotedblleft{}now\textquotedblright{}.

\begin{quote}
\emph{nǐ xiànzài máng ma} — Are you busy right now?
\end{quote}

Its characters come later; for now it is a word for the ear.

\begin{cousinweb}
\emph{xiànzài} is built from xiàn 'to appear' and zài 'to be present', both inherited from Middle Chinese, \emph{henH} and \emph{dzojX} (Guangyun, in Baxter's transcription).
\end{cousinweb}

\begin{grammarlens}[title={What you've built}]
One more word for times and partings.
\end{grammarlens}

\subsection*{Your turn}
\begin{itemize}
\raggedright
  \item \emph{Say it:} \emph{xiànzài}
  \item \emph{Say it:} \emph{xiànzài}, once more
  \item \emph{Say it:} \emph{nǐ xiànzài máng ma} — Are you busy right now?
  \item \emph{Recall:} say the Mandarin for today, then the Mandarin for yesterday, then say \emph{xiànzài} again
\end{itemize}

\begin{tcolorbox}[breakable,colback=teal!4,colframe=teal!35!black,title={Before you move on}]
What does xiànzài mean? (\textquotedblleft{}Now\textquotedblright{}.) Say it once more.
\end{tcolorbox}
\section[míngtiān]{míngtiān — tomorrow}

\label{lesson:ZH-C45-mingtian}



\begin{quote}
Before the new one: say the Mandarin for yesterday, then the Mandarin for now.
\end{quote}

\subsection*{You'll want to know: míngtiān}
\emph{míngtiān} — \textquotedblleft{}tomorrow\textquotedblright{}.

\begin{quote}
\emph{wǒ jīntiān hěn máng, míngtiān bù máng} — I'm busy today, but not tomorrow.
\end{quote}

Its characters come later; for now it is a word for the ear.

\begin{cousinweb}
\emph{míngtiān} is built from míng 'bright' and \zh{天} 'day', both inherited from Middle Chinese, \emph{mjaeng} and \emph{then} (Guangyun, in Baxter's transcription).
\end{cousinweb}

\begin{grammarlens}[title={What you've built}]
One more word for times and partings.
\end{grammarlens}

\subsection*{Your turn}
\begin{itemize}
\raggedright
  \item \emph{Say it:} \emph{míngtiān}
  \item \emph{Say it:} \emph{míngtiān}, once more
  \item \emph{Say it:} \emph{wǒ jīntiān hěn máng, míngtiān bù máng} — I'm busy today, but not tomorrow.
  \item \emph{Recall:} say the Mandarin for yesterday, then the Mandarin for now, then say \emph{míngtiān} again
\end{itemize}

\begin{tcolorbox}[breakable,colback=teal!4,colframe=teal!35!black,title={Before you move on}]
What does míngtiān mean? (\textquotedblleft{}Tomorrow\textquotedblright{}.) Say it once more.
\end{tcolorbox}
